\chapter{回归、分类与聚类的训练目标与评价指标}
\label{chap:objectives-metrics}
\begin{quote}\small
\textit{本章摘要。}一个模型需要先通过训练目标学习参数，再用评价指标检查它在新数据上的表现。训练目标规定什么样的错误会被惩罚，评价指标规定报告时关心什么结果。二者有时相同，有时不同。理解这种差别，才能解释为什么训练损失下降后，模型的实际表现仍可能不理想。本章分别讨论回归、分类和聚类，并说明每种任务中的常用目标、指标及其适用条件。
\end{quote}

训练集记为$\mathcal D_{\mathrm{train}}=\{(\bm x_i,y_i)\}_{i=1}^n$。模型用参数$\theta$给出预测$f_\theta(\bm x_i)$，再根据预测和标签之间的差异计算损失。常见的经验目标可以写成

\begin{equation}
 \hat{\theta}=\argmin_{\theta}\left[\frac1n\sum_{i=1}^n\ell(f_\theta(\bm x_i),y_i)+\Omega(\theta)\right].
 \label{eq:pipeline-general-objective}
\end{equation}
其中$n\geq1$是训练样本数，$\ell$是单个样本的损失，$\Omega$是可选的正则项。方括号中的量是训练经验目标。若$\Omega$不为零，它通常不是单个样本损失的平均值。验证集用于选择模型或训练设置，测试集用于最终评价。评价指标$M$在预先规定的数据划分和计算规则下衡量预测质量。若$M$不是$\ell$的同一个平均形式，训练目标与评价目标就存在差异，需要说明为何采用这种训练目标，以及如何检查最终关心的指标。

\section{回归：连续数值的预测}
设真实值为$y_i\in\R$，预测值为$\hat y_i=f_\theta(\bm x_i)$，误差记为$e_i=\hat y_i-y_i$。回归损失主要区别在于对大误差、方向性误差和异常值的处理方式。

\subsection{常用目标函数}
\paragraph{平方损失与均方误差。}
平方损失及均方误差（Mean Squared Error，MSE）为
\begin{equation}
 \ell_2(e)=\frac12e^2,\qquad
 \operatorname{MSE}=\frac1n\sum_{i=1}^ne_i^2.
 \label{eq:regression-mse}
\end{equation}
平方会放大大误差。例如，误差从$1$增至$3$时，单样本损失从$1/2$增至$9/2$，扩大为原来的$9$倍。在条件二阶矩有限时，平方损失的总体最优预测为
\begin{equation}
 \argmin_a\mathbb E[(Y-a)^2\mid X=\bm x]=\mathbb E[Y\mid X=\bm x].
 \label{eq:regression-conditional-mean}
\end{equation}
第一章式\eqref{eq:math-conditional-mse}已经将期望平方误差分解为条件方差与$(a-\E[Y\mid X=\bm x])^2$，因此最优预测是条件均值。这个结论不要求高斯噪声；高斯假设提供的是平方损失与负对数似然的对应。平方损失适合更重视大误差的任务，但对异常值敏感，其量纲是目标变量量纲的平方。

\paragraph{绝对损失与平均绝对误差。}
绝对损失及平均绝对误差（Mean Absolute Error，MAE）为
\begin{equation}
 \ell_1(e)=|e|,\qquad
 \operatorname{MAE}=\frac1n\sum_{i=1}^n|e_i|.
 \label{eq:regression-mae}
\end{equation}
在条件绝对一阶矩有限时，其总体最优预测是条件中位数；中位数可能不唯一：
\begin{equation}
 \argmin_a\mathbb E[|Y-a|\mid X=\bm x]=\operatorname{Median}(Y\mid X=\bm x).
 \label{eq:regression-conditional-median}
\end{equation}
绝对损失对异常值的增长是线性的，因而比平方损失稳健；代价是$e=0$处不可微，且当业务真正关心均值或极端大误差时可能不合适。

\paragraph{Huber损失。}
给定阈值$\delta>0$，Huber损失为
\begin{equation}
 \ell_\delta(e)=
 \begin{cases}
 \frac12e^2,&|e|\leq\delta,\\
 \delta(|e|-\frac12\delta),&|e|>\delta.
 \end{cases}
 \label{eq:regression-huber}
\end{equation}
小误差区域保留平方损失的平滑性，大误差区域改为线性增长。因此它在数值优化和异常值鲁棒性之间折中，但$\delta$会影响结果，通常需要通过验证集选择。

\paragraph{分位数损失。}
若任务关心预测上界、下界或非对称风险，可使用$\alpha$分位数损失：
\begin{equation}
 \ell_\alpha(e)=
 \begin{cases}
 (1-\alpha)e,&e\geq0,\\
 -\alpha e,&e<0,
 \end{cases}
 \qquad 0<\alpha<1.
 \label{eq:regression-quantile}
\end{equation}
本章定义$e=\hat y-y$，所以$e<0$表示低估，惩罚系数为$\alpha$；$e>0$表示高估，系数为$1-\alpha$。例如，取$\alpha=0.9$，同样大小的低估受到高估的$9$倍惩罚。在条件绝对一阶矩有限时，总体最优预测为条件$\alpha$分位数。若条件分布连续，在预测值$a$处的期望损失导数为$\Pr(Y\leq a\mid X=\bm x)-\alpha$，令其为零便得到分位数条件。$\alpha=0.5$时，损失是绝对损失的一半，二者最优解相同；$\alpha>0.5$适合更重视低估代价的任务。

\paragraph{相对误差和对数目标。}
平均绝对百分比误差（Mean Absolute Percentage Error，MAPE）可写为
\begin{equation}
 \operatorname{MAPE}=\frac{100\%}{n}\sum_{i=1}^n\left|\frac{e_i}{y_i}\right|.
 \label{eq:regression-mape}
\end{equation}
当$y_i$接近零时，MAPE会产生巨大权重，$y_i=0$时甚至无定义。对数变换或对称平均绝对百分比误差可以缓解部分问题，但也会改变误差的含义。只有在相对误差确实比绝对误差更接近业务代价时，才应使用这类指标。

\subsection{回归评价指标及差异}
\paragraph{误差大小。}
均方根误差（Root Mean Squared Error，RMSE）与前面的MAE都可以在独立评价集上计算：
\begin{equation}
 \operatorname{RMSE}=\sqrt{\frac1n\sum_{i=1}^ne_i^2},
 \qquad
 \operatorname{MAE}=\frac1n\sum_{i=1}^n|e_i|.
 \label{eq:regression-rmse-mae}
\end{equation}
RMSE与平方损失单调对应，且恢复到原目标的量纲；它对少量大错误更加敏感。MAE表示平均绝对偏差，解释直接且对异常值更稳健。二者没有普遍的优劣，差异来自错误代价函数不同。若RMSE明显大于MAE，通常说明误差分布中存在少数较大的偏差，但这不是异常值的唯一证据。

\paragraph{解释方差而非误差尺度。}
决定系数$R^2$为
\begin{equation}
 R^2=1-\frac{\sum_i(y_i-\hat y_i)^2}{\sum_i(y_i-\bar y)^2},
 \qquad \bar y=\frac1n\sum_i y_i.
 \label{eq:regression-r2}
\end{equation}
它比较模型与“始终预测测试集均值”的基线，$R^2=1$表示完美拟合，$R^2=0$表示平方误差与该基线相同，$R^2<0$表示平方误差更大。分母必须为正；当评价集中的标签全部相同时，这一公式无定义。$R^2$不具有目标变量的原始单位，且在不同数据集方差不同的情况下不宜直接比较；因此不能用它替代MAE或RMSE。

\paragraph{概率预测的评价。}
若模型输出预测分布$p_\theta(y\mid\bm x)$，可使用负对数似然
\begin{equation}
 \operatorname{NLL}=-\frac1n\sum_{i=1}^n\log p_\theta(y_i\mid\bm x_i).
 \label{eq:regression-nll}
\end{equation}
NLL同时评价预测中心和不确定性：预测过窄会因真实值落在低密度区域而受到重罚，预测过宽则会因密度分散而付出代价。仅报告均值的MAE或RMSE无法判断置信区间是否可靠。概率预测还应检查区间覆盖率与区间宽度，不能只看平均误差。

\section{分类：离散标签与概率决策}
设$y_i\in\{0,1\}$，模型输出正类概率$p_i=\Pr(Y=1\mid\bm x_i)$。分类中必须区分三种对象：概率预测、排序分数和阈值后的类别决策。它们对应的损失和评价指标不同。

\subsection{常用目标函数}
\paragraph{交叉熵与对数损失。}
二分类交叉熵为
\begin{equation}
 \ell_{\mathrm{CE}}(p,y)=-y\log p-(1-y)\log(1-p).
 \label{eq:classification-cross-entropy}
\end{equation}
若真实条件概率为$q=\Pr(Y=1\mid X=\bm x)$，预测为$p$，则期望损失
\begin{equation}
 q(-\log p)+(1-q)(-\log(1-p))
 =H(q)+\operatorname{KL}(q\|p)
 \label{eq:classification-cross-entropy-kl}
\end{equation}
其中$H(q)$是真实二分类分布的熵，$\operatorname{KL}(q\|p)$比较真实分布与预测分布。期望损失在$p=q$时达到最小。因此交叉熵的总体目标是正确的条件概率，但有限数据、受限模型和未充分优化都可能使预测偏离这一目标，实际概率仍需检查校准。对极少数类，普通交叉熵的总和可能主要由多数类贡献，需要类别权重、重采样或其他代价调整；这些方法改变了优化目标，不能再直接解释为原始概率的无偏拟合。

\paragraph{加权交叉熵和焦点损失。}
加权交叉熵为
\begin{equation}
 \ell_{\mathrm{WCE}}(p,y)=-w_1y\log p-w_0(1-y)\log(1-p).
 \label{eq:classification-weighted-ce}
\end{equation}
它通过提高某类错误的训练代价处理类别不平衡。焦点损失进一步降低容易样本的权重，例如
\begin{equation}
 \ell_{\mathrm{focal}}(p,y)=-(1-p_t)^\gamma\log p_t,
 \qquad p_t=yp+(1-y)(1-p),\quad\gamma\geq0.
 \label{eq:classification-focal}
\end{equation}
这类损失有助于把学习集中到困难样本，但概率值可能需要重新校准，不能默认仍具有良好的概率解释。

\paragraph{直接优化分类错误的困难。}
准确率对应零一损失
\begin{equation}
 \ell_{01}(\hat y,y)=\mathbb{1}\{\hat y\neq y\}.
 \end{equation}
它直接计数分类错误，但阈值决策使损失在参数空间的多数区域保持不变，在类别切换处发生跳变。普通梯度方法因而无法获得有效的局部改进方向。因此实践中常用交叉熵等可微替代损失，最后再用阈值产生类别。替代损失较小并不保证某个固定阈值下的$F_1$或业务成本已经最优。

\subsection{分类评价指标及差异}
\paragraph{从混淆矩阵确定统计对象。}
在异常检测或故障预警中，先把异常或故障定义为正类，把正常定义为负类。设模型给出分数$s_i$，分数越大表示越倾向于正类；只有分数经过概率建模或校准时，才把它解释为前文的$p_i$。给定阈值$t$，类别决策为$\hat y_i=\mathbb 1\{s_i\geq t\}$。固定阈值后，真正例（True Positive，TP）是正确报出的异常，假正例（False Positive，FP）是被误报为异常的正常样本，假负例（False Negative，FN）是漏掉的异常，真负例（True Negative，TN）是正确判为正常的样本。以下用$TP,FP,FN,TN$表示对应的样本数量。

准确率（Accuracy）、精确率（Precision）和召回率（Recall）分别为
\begin{equation}
 \operatorname{Accuracy}=\frac{TP+TN}{TP+FP+FN+TN},\quad
 \operatorname{Precision}=\frac{TP}{TP+FP},\quad
 \operatorname{Recall}=\frac{TP}{TP+FN}.
 \label{eq:classification-basic-metrics}
\end{equation}
准确率统计全部样本中判对的比例；精确率回答“报出的样本有多少是真的”；召回率回答“真实异常找到了多少”。三者的分母不同，因此回答的问题也不同。在异常稀少时，始终预测正常就可能得到很高的准确率，却漏掉全部异常。

\paragraph{误报率、漏报率与特异度。}
告警系统需要分别检查正常样本受到多少干扰，以及异常样本有多少没有被发现。假正例率（False Positive Rate，FPR）在这里称为误报率，假负例率（False Negative Rate，FNR）称为漏报率：
\begin{equation}
 \operatorname{FPR}=\frac{FP}{FP+TN},\qquad
 \operatorname{FNR}=\frac{FN}{TP+FN}.
 \label{eq:classification-fpr-fnr}
\end{equation}
误报率以全部真实正常样本为分母，漏报率以全部真实异常样本为分母。真正例率（True Positive Rate，TPR）就是召回率；真负例率（True Negative Rate，TNR）也称为特异度（Specificity）：
\begin{equation}
 \operatorname{TPR}=\operatorname{Recall}=1-\operatorname{FNR},\qquad
 \operatorname{TNR}=\frac{TN}{FP+TN}=1-\operatorname{FPR}.
 \label{eq:classification-tpr-tnr}
\end{equation}
这两组互补关系分别描述真实异常和真实正常样本的去向。

“误报率”在实际报告中有时还被用来表示$FP/(TP+FP)$。这个比值以全部告警为分母，称为假发现率（False Discovery Rate，FDR），等于$1-\operatorname{Precision}$。FPR回答“正常样本中有多少被误报”，FDR回答“告警中有多少是假的”。报告时应同时写出名称和计算式，避免用同一个词指代不同的量。

例如，10000个检测样本中有100个异常。某个阈值下，模型检出80个异常，并把99个正常样本判为异常，则$TP=80,FN=20,FP=99,TN=9801$。此时
\begin{equation}
 \begin{aligned}
 \operatorname{Accuracy}&=\frac{80+9801}{10000}=98.81\%,\\
 \operatorname{FPR}&=\frac{99}{9900}=1\%,\qquad
 \operatorname{FNR}=\frac{20}{100}=20\%,\\
 \operatorname{Precision}&=\frac{80}{179}\approx44.69\%,\qquad
 \operatorname{FDR}=\frac{99}{179}\approx55.31\%.
 \end{aligned}
 \label{eq:classification-alarm-example}
\end{equation}
正常样本中的误报比例只有$1\%$，但179次告警中有99次是假告警。正常样本数量远多于异常样本，少量正常样本被误报就可能占据大部分告警。同时，$98.81\%$的准确率没有反映出五分之一的异常被漏掉。

\paragraph{类别不平衡时的平衡准确率。}
普通准确率也可以写成各类召回率的加权平均。令测试集正例比例为$\hat\pi=(TP+FN)/n$，则
\begin{equation}
 \operatorname{Accuracy}=\hat\pi\operatorname{TPR}
 +(1-\hat\pi)\operatorname{TNR}.
 \label{eq:classification-accuracy-prevalence}
\end{equation}
当$\hat\pi$很小时，正常类的识别结果占据主要权重。若希望两类具有同等评价权重，可使用平衡准确率（Balanced Accuracy，BAcc）：
\begin{equation}
 \operatorname{BAcc}=\frac{\operatorname{TPR}+\operatorname{TNR}}2
 =1-\frac{\operatorname{FNR}+\operatorname{FPR}}2.
 \label{eq:classification-balanced-accuracy}
\end{equation}
它先分别计算两类的召回率，再等权平均。上例中$\operatorname{BAcc}=(0.80+0.99)/2=89.5\%$。始终预测正常的模型虽然有$99\%$的准确率，但$\operatorname{TPR}=0,\operatorname{TNR}=1$，所以BAcc只有$50\%$。因此，BAcc能揭示只依赖多数类获得的高准确率。

二分类BAcc取值在$[0,1]$；完美分类时为1，预测与标签独立时的总体值为$0.5$，也可能低于$0.5$。这里采用未进行机会水平校正的定义。多分类时，BAcc通常定义为各类别召回率的宏平均，即先按类别计算，再等权平均。它消除了类别数量在平均权重中的直接影响，但没有计入误报与漏报的实际代价，也不能说明告警中有多少是真的。比较模型时，仍应报告各类召回率、样本数和精确率。若某类在测试集中缺失，其召回率及相应的BAcc没有上述定义下的有效值。

\paragraph{同时检查精确率与召回率。}
若主要关心告警是否可信以及异常是否被找全，可用$F_1$概括精确率与召回率。它是二者的调和平均：
\begin{equation}
 F_1=\frac{2\,\operatorname{Precision}\,\operatorname{Recall}}
 {\operatorname{Precision}+\operatorname{Recall}}
 =\frac{2TP}{2TP+FP+FN}.
 \label{eq:classification-f1}
\end{equation}
调和平均会被较小的一项限制。上例中$F_1=160/279\approx0.573$，较低的精确率限制了结果。$F_1$没有使用$TN$，与同时检查两类召回率的BAcc侧重点不同。若希望调整精确率与召回率的相对权重，可使用
\begin{equation}
 F_\beta=\frac{(1+\beta^2)\operatorname{Precision}\operatorname{Recall}}
 {\beta^2\operatorname{Precision}+\operatorname{Recall}},\qquad\beta>0.
 \label{eq:classification-fbeta}
\end{equation}
$\beta>1$时更重视召回率，$\beta<1$时更重视精确率；$\beta$不是误报与漏报成本的直接比值。成本已知时，应直接计算决策成本。以上指标均须检查分母：没有真实正例时召回率无定义，没有正类预测时精确率无定义。若软件将无定义项置零或排除，应说明约定，不能把这种处理当作模型性能的证据。

\paragraph{从一个阈值到ROC曲线。}
降低阈值会让更多样本被判为异常。在同一测试集上，$TP$和$FP$都不会减少，召回率上升的同时，误报率也可能上升。接收者操作特征曲线（Receiver Operating Characteristic Curve，ROC曲线）把每个阈值对应的$\operatorname{FPR}(t)$作为横坐标，把$\operatorname{TPR}(t)$作为纵坐标，展示这两者的取舍。阈值高于所有分数时，没有告警，对应$(0,0)$；阈值不高于最低分时，全部样本告警，对应$(1,1)$。左上角$(0,1)$表示没有误报且没有漏报。

用六个样本即可看清曲线的形成过程。设三个异常样本的分数为$0.9,0.7,0.4$，三个正常样本的分数为$0.8,0.6,0.2$。由高到低的标签顺序是“异常、正常、异常、正常、异常、正常”。阈值从大到小越过一个异常样本时，$TP$增加1，TPR增加$1/3$，曲线向上移动；越过一个正常样本时，$FP$增加1，FPR增加$1/3$，曲线向右移动。由此得到图\ref{fig:classification-roc-example}中的阶梯曲线。

\begin{figure}[htbp]
\centering
\begin{tikzpicture}[x=5.4cm,y=5.4cm,font=\small]
 \fill[blue!10] (0,0)--(0,0.333333)--(0.333333,0.333333)
 --(0.333333,0.666667)--(0.666667,0.666667)
 --(0.666667,1)--(1,1)--(1,0)--cycle;
 \draw[gray!40,thin] (0.333333,0)--(0.333333,1)
 (0.666667,0)--(0.666667,1) (1,0)--(1,1)
 (0,0.333333)--(1,0.333333) (0,0.666667)--(1,0.666667)
 (0,1)--(1,1);
 \draw[gray!75,dashed] (0,0)--(1,1);
 \draw[->] (0,0)--(1.10,0) node[right] {FPR};
 \draw[->] (0,0)--(0,1.10) node[above] {TPR};
 \foreach \v/\lab in {0/0,0.333333/{\frac13},0.666667/{\frac23},1/1} {
  \draw (\v,0)--(\v,-0.015) node[below] {$\lab$};
  \draw (0,\v)--(-0.015,\v) node[left] {$\lab$};
 }
 \draw[blue!70!black,very thick] (0,0)--(0,0.333333)
 --(0.333333,0.333333)--(0.333333,0.666667)
 --(0.666667,0.666667)--(0.666667,1)--(1,1);
 \fill[red!75!black] (0.333333,0.666667) circle (2pt);
 \node[anchor=west,font=\scriptsize] at (0.37,0.75) {$t=0.7$};
 \node[align=center,fill=white,inner sep=3pt] at (0.70,0.23)
 {阴影面积\\$\mathrm{AUC}=2/3$};
\end{tikzpicture}
\caption{六个样本的经验ROC曲线。降低阈值时，异常样本使曲线向上移动，正常样本使曲线向右移动。红点对应阈值$t=0.7$，此时$TP=2,FP=1,FN=1,TN=2$；虚线为TPR等于FPR的参考线。}
\label{fig:classification-roc-example}
\end{figure}

例如，$t=0.7$时，分数为$0.9,0.8,0.7$的样本触发告警，工作点为$(1/3,2/3)$。把阈值降到$0.4$，三个异常都被找出，但已有两个正常样本被误报，工作点变为$(2/3,1)$。ROC曲线保留了这些候选工作点，便于检查提高检出率需要增加多少误报。若多个样本分数相同，它们会在同一阈值一起进入告警集合；通常将相邻可实现工作点用直线连接，中间点并不一定对应某个确定阈值。

\paragraph{ROC-AUC的面积与排序含义。}
曲线下面积（Area Under the Curve，AUC）把整条ROC曲线概括为一个数。上例的三个水平区间宽度均为$1/3$，高度依次为$1/3,2/3,1$，所以
\begin{equation}
 \operatorname{ROC\text{-}AUC}
 =\frac13\times\frac13+\frac13\times\frac23+\frac13\times1
 =\frac23.
 \label{eq:classification-roc-area-example}
\end{equation}
竖直线段宽度为零，不增加面积。采用相邻工作点间的线性连接时，经验AUC可用梯形面积求和。

同一个数还有排序解释。设正例数为$n_+=TP+FN$，负例数为$n_-=FP+TN$，正、负例的分数分别记为$s_i^+,s_j^-$，则
\begin{equation}
 \widehat{\operatorname{ROC\text{-}AUC}}
 =\frac{1}{n_+n_-}\sum_{i=1}^{n_+}\sum_{j=1}^{n_-}
 \left[\mathbb 1\{s_i^+>s_j^-\}
 +\frac12\mathbb 1\{s_i^+=s_j^-\}\right].
 \label{eq:classification-auc-pairs}
\end{equation}
它逐一比较正负样本对：正例排在前面记1分，排在后面记0分，同分记半分，再取平均。上例共有$3\times3=9$对。分数为$0.9$的异常排在三个正常样本前面，分数为$0.7$的异常排在两个正常样本前面，分数为$0.4$的异常排在一个正常样本前面，共有$3+2+1=6$对顺序正确，因此AUC为$6/9=2/3$，与面积计算一致。

在总体分布上，AUC等于独立抽取一个正例和一个负例时，正例得分较高的概率，加上同分概率的一半。AUC为1表示全部正例排在全部负例之前；无区分能力的随机分数在总体上对应AUC为$0.5$，有限样本估计会围绕它波动。AUC低于$0.5$可能提示分数方向反了，但也可能来自抽样波动或分布变化。AUC为$0.5$也不能证明分数与标签独立。

ROC-AUC只依赖排序。对分数施加严格单调递增变换，排序与AUC保持不变，概率含义却可能改变。因此，高AUC不保证概率校准，也不保证某个固定阈值下误报少。完整AUC汇总了全部误报率范围；若告警系统只允许FPR低于$1\%$，应重点比较该范围内的曲线，并报告满足误报约束时的TPR。整体面积更大的模型，在这个局部范围内未必更好。测试集需要同时包含正例与负例；负例太少时，经验FPR每次跳变较大，也无法可靠评价极低误报率下的性能。

\paragraph{稀有异常下的PR曲线。}
ROC曲线分别以正例数和负例数归一化。若两类各自的分数分布不变，仅正例比例变化，则总体ROC曲线与AUC不变，但告警中真假样本的比例会变化。精确率--召回率曲线（Precision--Recall Curve，PR曲线）以召回率为横轴、精确率为纵轴，直接展示“找出更多异常时，告警有多可信”。因此，在正例稀少的任务中，应同时检查PR曲线。

PR曲线下面积（Area Under the Precision--Recall Curve，PR-AUC）也可概括阈值变化下的表现，但必须说明计算约定。常用的平均精确率（Average Precision，AP）按召回率的增量加权精确率：
\begin{equation}
 \operatorname{AP}=\sum_k(R_k-R_{k-1})P_k,
 \label{eq:classification-average-precision}
\end{equation}
其中$R_k,P_k$是按分数从高到低扫描阈值得到的第$k$个工作点的召回率和精确率，$R_0=0$。上面的六样本例子中，召回率增加时的精确率依次为$1,2/3,3/5$，故$\operatorname{AP}=(1+2/3+3/5)/3=34/45\approx0.756$。AP采用这种阶梯式加权，与对PR曲线作线性插值后求梯形面积一般不同。没有正类预测时精确率无定义，绘图软件可能添加约定端点，也应区分这些端点与实际工作点。

无区分能力的随机分数在总体上对应精确率等于正例比例的PR基线。因此，PR-AUC或AP的数值须结合正例比例解释，跨数据集比较时不能忽略基率差异。这些面积仍然不能替代部署阈值下的误报率、漏报率和告警量。

\paragraph{概率质量：Brier分数与校准。}
二分类概率的Brier分数为
\begin{equation}
 \operatorname{Brier}=\frac1n\sum_{i=1}^n(p_i-y_i)^2.
 \label{eq:classification-brier}
\end{equation}
它同时惩罚概率偏离和过度自信。校准要求在预测概率约为$p$的样本组中，正例比例也约为$p$。交叉熵更强烈惩罚接近零或一但实际相反的预测，Brier分数的尺度更直观。二者都评价概率，不能与只评价排序的AUC混为一谈。

\paragraph{阈值、基率与业务代价。}
设总体正例基率为$\pi=\Pr(Y=1)$，某个阈值下的真正例率为$u$、假正例率为$v$。随机样本既为正例又触发告警的概率为$\pi u$，既为负例又触发告警的概率为$(1-\pi)v$。两者相加是触发告警的概率，因而
\begin{equation}
 \operatorname{Precision}=\frac{\pi u}{\pi u+(1-\pi)v}.
 \label{eq:classification-prevalence}
\end{equation}
分母须大于零。取$u=0.8,v=0.01$，当$\pi=0.01$时，精确率约为$44.69\%$；当$\pi$降至$0.001$而$u,v$保持不变时，精确率降至约$7.41\%$。这解释了为什么一个具有较高检出率和较低误报率的模型，部署到异常更稀少的环境后，仍可能产生大量假告警。

阈值选择应对应实际约束。例如，可以在验证集上寻找满足$\operatorname{FPR}\leq1\%$且召回率最高的阈值，或在每日告警量不超过人工处理能力的条件下比较漏报。若每次误报和漏报的成本分别为$c_{\rm FP},c_{\rm FN}$，并假设判对不产生成本，则每个样本的经验平均成本为
\begin{equation}
 \widehat C(t)=\frac{c_{\rm FP}FP(t)+c_{\rm FN}FN(t)}n.
 \label{eq:classification-alarm-cost}
\end{equation}
这里$FP(t),FN(t)$是阈值$t$下的错误数量；验证集应代表部署的类别比例，或按部署比例重新加权。成本不对称时，最大化BAcc或$F_1$不一定会最小化该成本。阈值和模型均应在验证阶段确定，测试集只用于最终评价，不能反复用于调阈值。

预警任务还必须明确统计单位。上述混淆矩阵默认每个样本作一次判断；连续时间窗中，同一次故障可能产生多个正例和多次告警。事件级评价应预先规定故障前多长时间内的告警算成功预警、如何匹配告警与故障，以及如何合并重复告警，再统计事件检出率、事件漏报率和预警提前量。正常运行期间每设备日的假告警次数衡量处理负担，单位为“次/设备日”，与以正常窗口数为分母的FPR不同。报告窗口级指标时，应同时说明窗口长度、步长和标签规则，不能直接把窗口召回率解释为故障事件检出率。

\section{聚类：没有标签时如何定义好坏}
聚类给定无标签样本$\{\bm x_i\}_{i=1}^n$，寻找分组$C_1,\ldots,C_K$。与回归和分类不同，聚类通常没有唯一的真实标签，因此必须先说明“好”的含义：簇内相似、簇间分离，还是与外部任务标签一致。

\subsection{常用目标函数}
\paragraph{$K$均值目标。}
$K$均值最小化簇内平方和：
\begin{equation}
 J_{\mathrm{kmeans}}(C,\{\bm\mu_k\})=
 \sum_{k=1}^K\sum_{i\in C_k}\|\bm x_i-\bm\mu_k\|_2^2,
 \qquad
 \bm\mu_k=\frac1{|C_k|}\sum_{i\in C_k}\bm x_i.
 \label{eq:clustering-kmeans}
\end{equation}
固定分组时，均值是平方距离的最优中心；固定中心时，把样本分给最近中心不会增加目标。交替执行这两步不会增加目标。在妥善处理空簇和相同距离的条件下，算法可停止于分配与中心都不再改变的状态，但不保证全局最优。该目标偏好近似球形、尺度相近且大小相近的簇，对特征缩放、异常值和$K$的选择敏感。

\paragraph{$K$中心与$K$中位数。}
$K$中心最小化最大簇内距离，强调最坏样本；$K$中位数使用
\begin{equation}
 J_{\mathrm{kmedians}}=\sum_{k=1}^K\sum_{i\in C_k}\|\bm x_i-\bm m_k\|_1.
 \label{eq:clustering-kmedians}
\end{equation}
它比平方距离对异常值更稳健，但通常仍需要预先给定$K$，并且所得分组取决于距离度量和特征尺度。

\paragraph{概率混合模型。}
高斯混合模型假设
\begin{equation}
 p(\bm x)=\sum_{k=1}^K\pi_k\mathcal N(\bm x\mid\bm\mu_k,\Sigma_k),
 \qquad \pi_k\geq0,\quad\sum_k\pi_k=1.
 \label{eq:clustering-gmm}
\end{equation}
最大化对数似然可以得到软分配，即每个样本属于各簇的后验概率。协方差矩阵允许簇具有不同方向和形状，表达力高于$K$均值，但也增加了参数、局部最优和退化协方差的风险。成分数可以依据独立验证数据的对数似然比较，也可以考虑加入参数数量惩罚的赤池信息准则（Akaike Information Criterion，AIC）或贝叶斯信息准则（Bayesian Information Criterion，BIC）。这些准则各有统计假设，且混合模型的退化情形需要另行处理。较好的概率拟合不保证分组具有实际任务所需的含义。

\subsection{聚类评价指标及差异}
\paragraph{内部指标。}
轮廓系数对样本$i$定义为
\begin{equation}
 s_i=\frac{b_i-a_i}{\max(a_i,b_i)},
 \label{eq:clustering-silhouette}
\end{equation}
其中$a_i$是样本$i$与本簇其余样本的平均距离。先分别计算它到各个其他簇的平均距离，再取最小值作为$b_i$。$s_i\in[-1,1]$；接近$1$表示更接近本簇，接近$0$表示处于分组边界附近，负值表示平均而言更接近某个其他簇。至少需要两个簇，且应说明单样本簇及分母为零时的处理约定。它依赖距离和簇形状，对非球形结构未必合适。

卡林斯基--哈拉巴斯指标（Calinski--Harabasz index，CH）以簇间离差与簇内离差的比值评价分组，较大值通常表示分离更明显。戴维斯--布尔丁指标（Davies--Bouldin index，DB）比较簇内散度与簇间中心距离，较小值通常表示簇更紧凑且相互分离。这些内部指标没有使用外部标签，只能评价分组是否符合各自的几何偏好。

\paragraph{外部指标。}
若存在独立的参考标签，可用调整兰德指数（Adjusted Rand Index，ARI）或归一化互信息（Normalized Mutual Information，NMI）比较聚类与参考分组。ARI检查样本对在两种分组中是否同簇，并校正随机分组也会产生的一致性；NMI衡量两种分组共享的信息，并进行归一化。它们适合回答“聚类是否复现某个已知分组”，但参考标签本身未必是自然结构；若标签是由业务规则产生，外部指标评价的是与该规则的一致程度，而不是无监督结构的绝对真理。

\paragraph{稳定性与任务效用。}
对数据进行重采样、扰动或改变初始化后，若分组频繁变化，说明结构不稳定。可用不同运行得到的分组之间的ARI或样本共聚类概率评价稳定性。若聚类用于下游检索、压缩或维护分组，还应测量下游任务效果。此时内部指标高并不必然意味着任务效用高：一个几何上紧凑的簇可能没有帮助决策的语义。

\section{三类任务的统一比较}
三类任务的目标函数和评价指标都可以写成对样本误差或结构差异的聚合，但聚合对象不同。回归直接比较连续误差，分类既可以比较概率也可以比较阈值决策或排序，聚类则通常比较距离结构、概率拟合或外部一致性。

\begin{table}[htbp]
\centering
\caption{三类任务中目标函数与评价指标的侧重点。}
\label{tab:pipeline-objectives-metrics}
\begin{tabular}{m{0.15\textwidth}m{0.25\textwidth}m{0.25\textwidth}m{0.25\textwidth}}
\toprule
任务 & 典型目标函数 & 典型评价指标 & 主要限制 \\
\midrule
回归 & MSE、MAE、Huber、分位数损失、NLL & MAE、RMSE、$R^2$、NLL、覆盖率 & 大误差、异常值、相对误差和不确定性侧重点不同 \\
分类 & 交叉熵、加权交叉熵、焦点损失 & FPR、FNR、BAcc、Precision、Recall、$F_1$、ROC-AUC、AP、Brier & 基率、阈值、统计单位、排序与概率校准不可混淆 \\
聚类 & $K$均值、$K$中位数、混合模型对数似然 & Silhouette、CH、DB、ARI、NMI、稳定性 & 无标签时不存在唯一评价标准，距离和$K$决定结论 \\
\bottomrule
\end{tabular}
\end{table}

因此，指标选择应从决策损失倒推。回归应同时报告至少一种原始量纲误差，并在需要时报告分位数或概率指标；分类应同时说明阈值性能、排序性能和概率质量，尤其要报告正例比例；聚类应明确使用的是内部结构、外部标签还是下游效用。只报告一个最高分，无法说明模型究竟改善了哪一种错误。

\section*{本章小结}
目标函数回答“训练时如何惩罚错误”，评价指标回答“报告时把什么视为成功”。平方损失偏向条件均值和大误差，绝对损失偏向条件中位数和稳健性；交叉熵评价概率信息，ROC-AUC评价排序，FPR和FNR分别检查正常样本的误报与异常样本的漏报，BAcc等权评价两类召回率，$F_1$评价固定阈值下的精确率--召回率平衡；$K$均值等聚类目标偏好特定几何结构，内部指标、外部指标和稳定性指标回答不同问题。实际报告必须同时写清数据划分、目标函数、阈值或簇数、评价指标及其适用假设。
